% =====================================================================
% Appendix A: AI Use Log
% =====================================================================
\subsection{AI Use Log}
\label{app:A}

\textbf{Method.} Before prompting AI for each task, the group
formulated its own initial ideas and analytical points.
We then followed the course cycle of generating an AI output,
evaluating and correcting it, and reflecting on its usefulness
and limitations. All five original task prompts were run in
Claude (Anthropic; model \texttt{claude-opus-5-5};
Claude desktop app) on 5~October 2026.
They were run in the same conversation used for source research,
so the model had access to research notes.
Each prompt was preceded by the context paragraph reproduced below.
The original outputs are preserved unedited in Appendix~B.

The group compared these outputs with its initial reasoning
and checked claims against course theory and case evidence.
Corrections were made to the analysis rather than retrospectively
changing the raw outputs. Our own further review
subsequently identified an overstatement in our own system-price
correction and problems with the classification of platform sides.
These further corrections are recorded below.

\begin{promptbox}
Context (given before all five prompts): Nord Pool is a European
power exchange owned 66\% by Euronext and 34\% by TSO Holding,
which is owned by the Nordic and Baltic transmission system
operators (including Statnett). For a university case study
in digital economics, treat the Euronext--TSO alliance behind
Nord Pool as the unit of analysis, not a single firm.
Keep each answer to max 400 words.
\end{promptbox}

% ---------------------------------------------------------------------
\subsubsection*{A.1 Task 1: Business ecosystem}

\textbf{Tool:} Claude, \texttt{claude-opus-5-5},
5~October 2026.
\quad \textbf{Output:} Appendix~B.1.

\begin{promptbox}
Identify the direct and indirect stakeholders in Nord Pool's
business ecosystem, the key interactions, complementarities
and dependencies between them, and explain how the alliance
structure changes the roles, incentives and influence
of the main actors.
\end{promptbox}

\textbf{Claims checked.} The current owners of TSO Holding
(EPSO-G, Statnett, Svenska kraftnät) are correct according to
\citet{NPAbout}. The output also fixes an error in our own
context paragraph, which implied that all Nordic and Baltic
TSOs are owners. The exit of Fingrid and Energinet matches
\citet{EPSOG2022}, and ``around 400 firms'' matches
\citet{NP2024}. EPEX's entry in 2020 matches \citet{EPEX2020}.
We also kept two points from the output after checking them.
Its description of Euronext and the TSOs as bringing
complementary resources (commercial and technological
capabilities on one side; system knowledge, capacity data
and legitimacy on the other) is consistent with Euronext's
stated aims \citep{Euronext2020} and the TSOs' role in
price calculation \citep{NPPrice}; we grounded it in the
complementarity concept of \citet{Jacobides2018}
(Section~\ref{sec:t1}). Its point that the TSOs' influence
now runs through regulation and the shared coupling rather
than ownership is consistent with their single board seat
\citep{NPBoard} and the 2022 exits \citep{EPSOG2022}.

\textbf{Error 1: oversimplified dependency model.}
The output describes Nord Pool and the TSOs as a symmetric
``mutual dependency''. Applying the textbook's dependency
notation \citep[Fig.~4.3]{OverbyAudestad2021} exposes
the oversimplification. Nord Pool relies on TSO capacity inputs,
but under the multi-NEMO arrangement the TSOs depend on the
shared coupling rather than on Nord Pool alone
\citep{EPEX2020}. The claim that without Nord Pool's auction
there would be ``no efficient capacity allocation or congestion
income for the TSOs'' is also overstated,
since congestion income arises through the coupled market
\citep{EU2019943}.

Our initial revision introduced a further overstatement
by claiming that the Nordic system price does not depend
on TSO data. Nord Pool's documentation states that internal
Nordic capacities are set to infinity, but that flows between
the Nordics and surrounding markets from the area-price
calculation are included as import or export orders
\citep{NPPrice}. Removing internal congestion restrictions
therefore does not remove all dependence on capacity inputs.
\emph{Correction:} Sections~\ref{sec:t1} and~\ref{sec:t5}
now distinguish the two calculations without claiming
that the system price is independent of TSO data.
This correction also shows why our revisions required
further checking after the first AI evaluation.

\textbf{Error 2: unsupported claim about preferences.}
The output states that ``TSOs want transparent, low-cost
price formation''. It gives no evidence, and we found no
source stating the TSOs' preferences in these terms.
\emph{Correction:} we replaced attributed motives with
observable evidence: the TSOs hold one of six shareholder board seats
\citep{NPBoard}, two TSOs exited ownership
\citep{EPSOG2022}, and Euronext has made documented
commercial decisions on data and futures
\citep{ERR2022,Euronext2026}.

\textbf{Error 3: outdated futures wording and a misclassified owner.}
The output states that Euronext ``since 2026 also runs
Nordic power futures in the Nord Pool environment''.
This follows wording from before the launch: when the
acquisition closed, Euronext's CEO said the new platform
``will be operated in the NordPool environment with clearing
by Euronext Clearing'' \citep{FOW2025}. The March 2026
launch release is more precise: the futures trade on
Euronext's Optiq platform and clear on Euronext Clearing,
while Nord Pool provides the underlying system and spot
prices \citep{Euronext2026}. The distinction matters for
the alliance analysis, because it shows which partner runs
the futures market. The output also calls EPSO-G,
Statnett and Svenska kraftnät ``the TSOs'', but EPSO-G
is the Lithuanian state-owned group that owns the TSO
Litgrid, not a TSO itself \citep{EPSOG2022}.
Both errors first slipped into our own draft and were
caught when we checked the draft against the sources.
\emph{Correction:} Section~\ref{sec:t1},
Table~\ref{tab:stakeholders} and Figure~\ref{fig:ecosystem}.

\textbf{Missing stakeholder.} The output omits EEX,
whose Nordic system price futures use the system price
that Nord Pool calculates \citep{EEX2024}.
This illustrates the textbook's point that removing
or changing one actor affects others in the ecosystem
\citep[Sect.~4.1]{OverbyAudestad2021}.
We added EEX to Table~\ref{tab:stakeholders}.

\textbf{Still unverified.} The TSOs' actual objectives in
the alliance and the employers of four board members are not
public. We would need TSO Holding's shareholder agreement or
the owners' annual reports before treating any claim about
TSO motives as fact.

\textbf{Reflection.} The output was useful as a fast,
well-structured first list of stakeholders.
It failed where judgement was needed: it treated an
asymmetric relationship as symmetric and presented
assumed motives as facts. Further checking also showed
that our own first correction could introduce an error.

% ---------------------------------------------------------------------
\subsubsection*{A.2 Task 2: Network effects and path dependency}

\textbf{Tool:} Claude, \texttt{claude-opus-5-5},
5~October 2026.
\quad \textbf{Output:} Appendix~B.2.

\begin{promptbox}
Identify the network effects associated with Nord Pool's
platform. Does it become more valuable as more participants,
market areas, TSOs, data sources or complementary services
join? Do data accumulation, standards, sunk costs,
learning effects or switching costs create path dependency,
lock-in, tipping or winner-takes-all dynamics?
\end{promptbox}

\textbf{Claims checked.} The system price as the reference
for Nordic financial contracts \citep{NPPrice},
multi-NEMO in 2020 \citep{EPEX2020} and the 15-minute
MTU \citep{OTE2025} are correct.
The point that the coupling is shared by all exchanges
is consistent with how SDAC works.
We kept the output's overall conclusion that the Nordic
spot market had tipped to Nord Pool before 2020 and is now
a contestable incumbent rather than a winner-takes-all
monopoly. Until the multi-NEMO go-live on 3~June 2020,
the Nordic bidding zones were not open to more than one
exchange \citep{EPEX2020}, and EPEX's 54~TWh in 2024
\citep{EPEX2025} shows a foothold rather than a takeover
(Section~\ref{sec:t2}). The five effects in
Table~\ref{tab:ne} also follow the output's list;
we kept each one after checking its classification against
the textbook \citep[Sect.~9.2--9.5]{OverbyAudestad2021}
and adding a source for the mechanism.

\textbf{Error 1: missing key distinction.}
The output confuses network effects with economies of scale.
An exchange's low cost per additional trade is a supply-side advantage.
A network effect occurs when participation changes the value
of the service for other users
\citep[Def.~9.1]{OverbyAudestad2021}.
\emph{Correction:} Section~\ref{sec:t2}
opens with this distinction.

\textbf{Error 2: unsupported claims.}
The output asserts that each new bidding zone or interconnector
``increases the set of feasible trades and welfare''
without evidence, and that ``the marginal value of one more
large Nordic member is smaller'' without any basis.
\emph{Correction:} we support the market-integration
argument with \citet{Newbery2016} and dropped
the unsupported marginal-value claim.

\textbf{Error 3: generic argument.}
``Sunk costs and learning effects in market design reinforce
the Nordic model'' names no mechanism and no evidence.
\emph{Correction:} we kept both concepts, which the task
names, but stated their mechanisms: sunk investments
in IT integration and collateral arrangements,
and learning effects among staff trained on Nord Pool's
systems. Both act as switching costs
\citep[Def.~12.2]{OverbyAudestad2021}.
We added the de facto standard argument
\citep[Def.~15.1]{OverbyAudestad2021}
and \citet{Arthur1989} on lock-in from early advantages.
We also added the observation that regulation
\emph{lowers} switching costs, the electricity analogue
of number portability
\citep[Sect.~12.6]{OverbyAudestad2021}.

\textbf{Missing evidence.}
``EPEX has entered'' and ``Lock-in may be moving to the system
price, data and futures''
were not quantified. We added EPEX's 54~TWh of Nordic
day-ahead volume in 2024 \citep{EPEX2025}
and the open-interest figures for Euronext and EEX
\citep{Euronext2026,EEX2026}.
We also classified the effects in Table~\ref{tab:ne},
which the output did only partly.

\textbf{Still unverified.} We have EPEX's Nordic volume for
one year only, so we cannot tell whether its share is
growing. Multi-year volumes per exchange would be needed to
judge whether the spot market is tipping back.

\textbf{Reflection.} The list of effects was a usable
checklist. Its weakness was confident claims without
evidence and a missing conceptual distinction that
the course treats as fundamental.

% ---------------------------------------------------------------------
\subsubsection*{A.3 Task 3: Multisided platform and pricing}

\textbf{Tool:} Claude, \texttt{claude-opus-5-5},
5~October 2026.
\quad \textbf{Output:} Appendix~B.3.

\begin{promptbox}
Identify the distinct user groups (sides) of Nord Pool's
platform and the cross-side and same-side network effects
between them. Describe Nord Pool's pricing model for
each side: who pays, how much, and who is subsidised.
\end{promptbox}

\textbf{Claims checked.}
The principal trading fees quoted by the output match
the official 2026 fee schedule \citep{NPFees2026}.
The €3,500 retailer licence matches \citet{NPData}.
The €600 historical-data price matches the 2022 report
\citep{ERR2022}, but this does not establish that
it remains the price in 2026.
The output did not invent these quoted prices,
although their dates and scope still required checking.
We also kept two points from the output. That buyers and
sellers pay the same fees, so there is no free side,
is confirmed by the published tariff \citep{NPFees2026}.
The figure that fees are about 0.15\% of the average price
also comes from the output; we recalculated it as
€0.055/MWh (€0.040 trading + €0.015 clearing) divided by
the 2024 average system price of €36.06/MWh \citep{NP2024},
which gives 0.15\%.

\textbf{Error 1: missing theoretical test.}
The output assumes that Nord Pool is an MSP without
testing the definition
\citep[Def.~10.1]{OverbyAudestad2021}.
Its central-counterparty role also requires distinguishing
mediation from resale.
\emph{Correction:} Section~\ref{sec:t3}
applies the control-over-terms criterion of
\citet{HagiuWright2015}
to the core buyer--seller interaction.

\textbf{Error 2: overlapping groups treated as sides.}
Small participants are buyers or sellers using a different
tariff, not an additional side merely because of that tariff.
Futures traders are affiliated with complementary markets,
and buying a data product does not by itself demonstrate
direct interaction with another affiliated user group.
Our first version repeated this classification problem.
\emph{Correction:} Section~\ref{sec:t3}
and Figure~\ref{fig:msp} now distinguish the two core
trading sides from tariff segments, ancillary services
and complementary markets.
Figure~\ref{fig:srm} preserves these actors as stakeholders
without treating them all as platform sides.

\textbf{Error 3: unsupported subsidy interpretation.}
The claim that data customers and mid-sized members
are the main ``money side'' is not established
by the published tariff.
Lower fixed fees and caps may encourage participation,
but do not by themselves identify who finances whom.
\emph{Correction:} we dropped data customers as a money
side and kept the rest as an explicit, labelled assumption
rather than a finding. Section~\ref{sec:t3} states the
assumption, defends it from the fee structure and notes
that the tariff alone cannot show who finances whom,
which is what the task asks for when pricing is not fully
disclosed. We interpret the differentiated pricing using
\citet{RochetTirole2003} and \citet{Armstrong2006}.
The Small Participant annual floor is included
in the revised platform figure.

\textbf{Error 4: oversimplified TSO costs.}
``TSOs do not pay'' overlooks the sharing of coupling
costs under CACM \citep{CACM2015}.
\emph{Correction:} we distinguish the absence of trading
fees in the capacity-provider role from contributions
to the coupling infrastructure.

\textbf{Still unverified.} Who finances whom cannot be
confirmed without Nord Pool's revenue broken down by
customer group, which it does not publish. Confirming
whether the €600 historical-data price still applies in 2026
would require Nord Pool's current price list.

\textbf{Reflection.}
The output reproduced several sourced prices accurately
but blurred the distinction between platform sides
and other customer or stakeholder groups.
Revising the output required applying the MSP definition
consistently to the text, pricing table and figures.

% ---------------------------------------------------------------------
\subsubsection*{A.4 Task 4: Cybersecurity economics}

\textbf{Tool:} Claude, \texttt{claude-opus-5-5},
5~October 2026.
\quad \textbf{Output:} Appendix~B.4.

\begin{promptbox}
How should Euronext and the TSOs allocate cybersecurity
responsibility and investment for Nord Pool's platform
between them? Would hosting and operating the platform
domestically in the Nordic region make it more secure?
\end{promptbox}

\textbf{Claims checked.}
The output does not invent probabilities, costs
or return-on-investment figures, and it treats security
as a jointly provided good, consistent with
\citet{Kianpour2022}.
NIS2's listing of NEMOs \citep{NIS2}
and the start date of the Norwegian Digital Security Act
\citep{Regjeringen2025} are correct.
Its argument about private incentives identifies
a security-externality mechanism
\citep{AndersonMoore2006}, although it does not measure
the alliance's actual security investment.
Two further parts of the output appear accurate, and we
kept them. Its division of the attack surface (core trading
and clearing systems with Nord Pool, capacity data with the
TSOs, trading systems and credentials with the members,
incident response shared within the coupling) fits the
provision structures of \citet{Hirshleifer1983} and is
consistent with the decoupling incidents
\citep{NPIncident2019,SDAC2021}; Section~\ref{sec:t4}
builds on it. Its answer on domestic hosting is also kept:
the weak points are interfaces, credentials and capacity
data rather than the data-centre location, and Euronext is
EU-based. We added the point about extraterritorial access
by non-EU providers. Public sources do not state where
Nord Pool's systems are hosted, so this answer rests on the
attack-surface reasoning rather than on evidence about the
location.

\textbf{Error 1: instruments implied to exist without evidence.}
The output lists ``SLAs with penalties between Nord Pool
and the TSOs, \dots\ joint exercises and cyber insurance''
without indicating whether the alliance actually uses them.
We found no public evidence confirming these arrangements.
\emph{Correction:} Section~\ref{sec:t4}
states explicitly where the public record is silent.

\textbf{Error 2: flawed allocation rule.}
The output recommends allocating investment
``in proportion to control over each attack surface''.
Under weakest-link provision, the least-protected
contribution determines the level of protection
\citep{Hirshleifer1983}, while losses can also fall
on actors who control no part of the platform.
Allocating investment by control alone does not
necessarily resolve this externality.
\emph{Correction:} our cost and loss map
(Table~\ref{tab:cyber}) separates control from loss
and considers instruments for reallocating risk.

\textbf{Error 3: unsupported assumption.}
``Euronext can bring group-level security operations,
threat intelligence and resilience standards''
is plausible, but no source describes how Nord Pool's
security is integrated with Euronext.
\emph{Correction:} Section~\ref{sec:t4}
labels it explicitly as our assumption.

\textbf{Missing links.}
The task requires linking the loss magnitude
to the network effects in Task~2
and the pricing incentives in Task~3.
The output does neither and includes no cost and loss table.
We added these connections and the table.
Following the revised platform classification,
Small Participants are discussed as a tariff segment
within the trading sides rather than a separate
subsidised platform side.

\textbf{Still unverified.} How security operations are split
between Nord Pool and Euronext, and where the systems are
hosted, would require disclosures from the companies
themselves; our cost and loss levels therefore remain
illustrative.

\textbf{Reflection.}
The output was cautious and well structured,
which made it a useful starting point.
It was weakest where it moved from description
to recommendation, and it did not connect security
to the rest of the analysis.

% ---------------------------------------------------------------------
\subsubsection*{A.5 Task 5: Business Model Canvas and SRM}

\textbf{Tool:} Claude, \texttt{claude-opus-5-5},
5~October 2026.
\quad \textbf{Output:} Appendix~B.5.

\begin{promptbox}
Create a Business Model Canvas (all nine blocks)
for the Nord Pool alliance, and a Stakeholder Relationship
Model showing the key stakeholders, their relationships,
dependencies, and positive or negative network effects.
\end{promptbox}

\textbf{Claims checked.}
The blocks are broadly consistent with Nord Pool's
own description of its services \citep{NPAbout}
and with the textbook's block definitions
\citep[Table~19.1]{OverbyAudestad2021}.
The value proposition in Section~\ref{sec:t5} (liquid,
neutral price formation with guaranteed settlement) is kept
from the output; it matches Nord Pool's description of its
markets and its role as central counterparty
\citep{NPAbout,NPClearing}.

\textbf{Error 1: unsupported revenue claim.}
``Futures (via Euronext)'' is listed as a Nord Pool
revenue stream without evidence.
The futures trade on Euronext's Optiq platform
and clear on Euronext Clearing, and the press release
states that ``Nord Pool provides the underlying system
and spot prices'' \citep{Euronext2026}.
How the futures revenue is split between Euronext
and Nord Pool is not public.
Asking \emph{whose} revenue it is exposes the gap,
but the evidence does not support the opposite claim
either, that all of it goes to Euronext.
\emph{Correction:} Section~\ref{sec:t5}
and Figure~\ref{fig:bmc}
state that the split is not public.

\textbf{Error 2: single-firm framing of partners.}
``TSOs'' are treated as one partner group.
This hides an important governance distinction:
some TSOs participate in ownership and others
are input suppliers without an ownership stake
\citep{EPSOG2022}.
The output also lists EPSO-G among the TSOs
(see A.1).
\emph{Correction:} we distinguish owner
and non-owner TSOs.

\textbf{Error 3: missing cost item.}
The cost structure omits the shared costs
of the coupling, which CACM allocates
between NEMOs and TSOs \citep{CACM2015}.
\emph{Correction:} these costs are included
in the revised canvas.

\textbf{Unsupported segment.}
``TSOs (auction services)'' appears as a customer
segment without evidence for the specific claim.
We could not verify it and removed it.

\textbf{SRM notation and platform boundary.}
The initial SRM was provided in text form.
We drew Figure~\ref{fig:srm} using the textbook
notation and the explicitly marked extension (+/0)
for the informational benefit of trading activity
for data customers.
Our own subsequent review also identified that our figures
treated Small Participants as a distinct side
even though they are a tariff segment
within buyers and sellers.
We removed the separate cross-side-effect arrow
for this segment.
Futures traders and data customers remain
in the SRM as stakeholders in complementary markets
and ancillary services.
Being included in an SRM does not imply being a side
of the core MSP
\citep[Def.~10.1 and Sect.~19.3]{OverbyAudestad2021}.

\textbf{Still unverified.} The split of futures revenue
between Euronext and Nord Pool, and the size of each cost
item, would require segment figures from Euronext's or Nord
Pool's annual reports.

\textbf{Reflection.}
The canvas was a complete first checklist.
It reproduced the weakness of writing a canvas
as if one firm owned all costs, revenues and resources.
The revised SRM also distinguishes the wider stakeholder
network from the narrower set of core platform sides.

% ---------------------------------------------------------------------
\subsubsection*{A.6 Other AI use: source research, tables, figures and LaTeX conversion}

\textbf{Tool:} Claude (Anthropic), Claude desktop app,
\texttt{claude-opus-5-5}, October 2026.

\textbf{Use.} Source research, in the same conversation as
the prompts in A.1--A.5. Claude helped find candidate case
sources (company pages, press releases, regulations) and
peer-reviewed articles. Every source cited in this report was
opened and checked by the group, and no AI tool is cited as
evidence for any fact.

\textbf{Tool:} Claude (Anthropic), Claude desktop app,
\texttt{claude-opus-5-5}, 5--6~October 2026.

\textbf{Use.} The group gave Claude its own sketches of how
the tables and figures should look, and Claude produced the
tables and figures in the report from these sketches.
Later revisions of Figures~\ref{fig:msp} and~\ref{fig:srm}
were converted by ChatGPT (see below). The content of the
tables and figures (actors, classifications, levels and
prices) comes from the group's analysis and was checked by
the group against the main text and sources.

\textbf{Tool:} ChatGPT (OpenAI), model GPT 6.1 Sol.

\textbf{Use.} The group wrote its later revisions in plain text:
the corrected system-price dependency,
the revised classification of platform sides,
and the updated declaration, Appendix~A and reflection text.
ChatGPT was used only to convert this text into LaTeX code,
including the TikZ code for Figures~\ref{fig:msp}
and~\ref{fig:srm}, and to commit it to the project repository.
The content and wording of these revisions are the group's own.
The original five task outputs in Appendix~B
remain unchanged.

\textbf{Tool:} Claude (Anthropic), Claude Desktop, Opus 5.5, 7~October 2026.

\textbf{Use.} The group worked on the text in a Word
copy of the report. Claude was used only
to convert that wording back into the LaTeX files, keeping the
existing citations, cross-references, tables and figures.

% ---------------------------------------------------------------------
